\section{Three connected research questions}
Generating a deliverable, and the documentation meant to vouch for it, is now nearly free while independent verification remains costly. When anyone can fabricate the evidence, the question is not what to disclose but which rules make truthful disclosure stable. The stakeholders are plugin-marketplace developers, the platforms setting the rules, and the users and journalists who act on a claim.

Three foundations frame the problem. Akerlof~\cite{akerlof1970} shows unverifiable quality drives good sellers out of a market. Shalvi et al.~\cite{shalvi2012} show dishonesty grows with available justifications. Buneman et al.~\cite{buneman2001} define provenance as structured lineage, but the producer writes the record, so the trace itself can be forged.

\textbf{Q1---Economics:} Which mechanisms---a fine $pF$ on detected fakes, or a market premium $\Delta$ for verifiable provenance---make honesty a Nash equilibrium when trace files are free to generate and to fake?\par
\textbf{Q2---Computation:} What trace format and verifier make a per-claim declaration checkable at near-zero cost, without trusting the producer's own record?\par
\textbf{Q3---Behavior:} Does traceability shrink the justification space, or does effortless checking produce moral licensing that increases dishonesty?

The three are jointly necessary: a rule requiring trace files is empty unless a fake can be cheaply detected, and an enforceable rule may still fail if checking licenses inattention. Figure~\ref{fig:teaser} connects them.

\section{Answer to Q1: incentives and intellectual lineage}
\textbf{Hypothesis.} Either mechanism suffices. A penalty $pF$ falls only on a detected fake; a premium $\Delta$ is earned only by a verifiable trace file. Honesty is the unique best response when $\Delta + pF$ exceeds the fake's cost advantage.

\textbf{Model.} A developer chooses an action $a_i \in \{\text{honest, none, fake}\}$ and verifier enrolment $v_i \in \{0,1\}$. A platform commits to $(m,p,F)$ for mandatory attachment, spot-check rate and fine; the market sets $\Delta$. Payoff is market benefit plus $\Delta$ if honest, minus cost, plus stability if enrolled, minus $pF$ if faking. A fake earns the same base benefit as an honest deliverable, since buyers cannot observe truth before purchase, while saving cost. Information is incomplete. The solution concept is Nash equilibrium in pure strategies~\cite{nash1950}, with $(m,p,F)$ exogenous; because the platform chooses the rules developers best-respond to, this is a mechanism-design problem.

\begin{table}[h]\caption{Strategy matrix with no rules ($p=F=0$). Left: no premium, so faking wins. Right: a premium, so honesty wins with no enforcement.}\label{tab:matrix}
\renewcommand{\arraystretch}{1.18}
\begin{tabular}{@{}lrr@{}}\toprule
 & \textbf{$\Delta=0$} & \textbf{$\Delta=30$} \\\midrule
A \quad honest trace & 135 & \textbf{165}$^{\star}$ \\
B \quad no trace     & 120 & 120 \\
C \quad faked trace  & \textbf{160}$^{\star}$ & 160 \\\bottomrule
\end{tabular}\end{table}

Without a premium, faking wins by $25$---exactly the cost it avoids---and only a fine can deter it: $p^{*}=0.5$ at $F=50$. With a premium, honesty wins by being worth more. The mechanisms substitute: $\Delta=15$ halves the required spot-check rate to $0.2$, and $\Delta=25$ removes the need to fine at all.

\textbf{Unresolved.} Akerlof fixes buyers' beliefs rather than asking how a costly signal forms them. Spence~\cite{spence1973} supplies that logic: a premium exists only if the signal separates types.

\section{Answer to Q2: method and evaluation}
\textbf{Method.} A browser-run $3\times2$ game on a Hugging Face Space, client-side HTML, CSS and JavaScript; the payoff law is in Appendix~\ref{app:tech}. The algorithm enumerates the six cells and finds the best response.

\textbf{Baseline correction.} The first deployed version summed a cost, a benefit and a stability index. Cost and benefit have opposite signs, so summing them scored honesty as better on both; the sweep returned $p^{*}=0$, meaning faking was never attractive and the question could not be posed. The corrected law enters cost negatively and gives the fake the same benefit as an honest deliverable. See Appendix~\ref{app:develop}.

\textbf{Baseline, modification, metric, check.} Baseline is the status-quo cell (net $120$); the modification is mandatory attachment plus $p$ and $F$; the metric is the threshold $p^{*}$ at which honesty is uniquely best. Six checks pass: no rules $\Rightarrow$ fake wins; enforcement alone makes honesty unique; a premium alone ($\Delta=30$, $p=F=0$) also does; the status quo is unchanged; enrolment is weakly dominant; six cells.

\textbf{Result.} At $F=50$, $p^{*}=0.5000$; the required $p$ falls as $F$ rises ($F=100\Rightarrow0.25$; $F=400\Rightarrow0.0625$), and below $F=25$ no $p\le1$ suffices. All results are \emph{simulated}; the Colab notebook reproduces them from one payoff law.

\section{Answer to Q3: behavior and competing explanations}
\textbf{Competing explanations.} Justification: dishonesty grows with available justifications~\cite{shalvi2012}, so traceability should shrink that space. Moral licensing: a visible attestation lets producer and verifier treat the check as discharging responsibility, so the producer shades a claim the verifier no longer examines. They predict opposite effects from the same manipulation.

\textbf{Discriminating evidence.} They separate on verifier attention. Justification predicts less faking with scrutiny unchanged; licensing predicts reduced scrutiny and a fake rate falling less than the attachment rate rises. Measuring both producer choice and check depth discriminates; measuring the fake rate alone cannot.

\textbf{Consequence and limits.} If licensing dominates, verification must not be a binary pass: it should report which claims were recomputed, and $pF$ must be paired with something keeping scrutiny positive. A simulated table cannot establish human behavior; the next test crosses trace-file presence with check depth.

\section{Advanced development and future directions}
\textbf{Foundation.} Akerlof's quality-uncertainty analysis~\cite{akerlof1970}, recognized by the 2001 Nobel Memorial Prize in Economic Sciences~\cite{nobel2001}, establishes that unverifiable quality destroys markets and that institutions making quality observable restore trade. The enduring question is trust under asymmetric information. Spence~\cite{spence1973} supplies the complementary logic: a costly credible signal separates quality only if the receiver can read it.

\textbf{The boundary and the advance.} Build-layer provenance~\cite{slsa} already exists; unresolved is the claim layer, where per-claim sources and input--output declarations are recomputed at near-zero cost although the producer writes the record and can forge it. Treating $\Delta$ as a constant merely relocates the assumption: a premium exists only if publishing a trace is informative about type, which is Spence's single-crossing condition. If fakers can publish traces as cheaply as honest producers the signal pools, $\Delta$ falls to zero, and enforcement is again the only mechanism. Making $\Delta$ endogenous is the next study; the certification-design result of M{\"a}kimattila et al.~\cite{makimattila2025} is the natural instrument, since there a sender's \emph{choice} among offered tests is itself the signal---but in that model the receiver obtains zero surplus, so a premium does not reach honest producers automatically.

\textbf{The verifier is built for machines.} It returns pass or fail, which a non-specialist cannot interpret; letting an ordinary user judge a firm's claim, or its data-handling practice, is a problem of cognitive rather than computational cost, where the objective is accurate enough and fast enough rather than perfect.

\begin{table}[h]\caption{From laureate foundations to a new interdisciplinary contribution.}
\begin{tabularx}{\columnwidth}{@{}p{1.5cm}Y@{}}\toprule
\textbf{Horizon}&\textbf{Milestone and evidence}\\\midrule
Foundation&Akerlof (1970): unverifiable quality collapses markets. Spence (1973): a credible signal needs single crossing.\\
PS1&A $3\times2$ game; two mechanisms located, a fine $pF$ and a premium $\Delta$, with $\Delta+pF>25$ required.\\
Next study&Endogenise $\Delta$ in buyer beliefs: does trace publication separate types or pool?\\
Longer term&A verifier a non-specialist can use, reporting what was checked, at what accuracy and cost.\\
2056&Verifiable honesty that is competitively rewarded, not merely enforced, and readable by any user.\\\bottomrule
\end{tabularx}\end{table}

\noindent\textbf{Expected 2056 Nobel Prize and Turing Award contribution statement.}
By 2056, I aim to contribute a governed trust infrastructure for an era of free generation, in which verifiable provenance is competitively rewarded rather than merely enforced, and in which any ordinary user can check a firm's claims and data-handling practices without technical training. Building on Akerlof's quality-uncertainty foundation, Spence's signalling logic and machine-checkable provenance, I will integrate economics, computer science and behavioral science so that near-zero-cost trace files, a verifier designed for non-specialists, and platform rules that make honest publication an equilibrium together let anyone inspect a claim in under a minute. Progress will be shown by an endogenous premium that separates honest from faking producers, a verifier whose accuracy and cost are documented, a discriminating behavioral test, and a rule design adopted on a live platform.

\subsection*{Open Science Statement}
\textbf{GitHub:} \GitHubLink\\
\textbf{Google Colab:} \ColabLink\par
Released: a browser-run simulator and a Colab notebook that re-derives the payoff table and runs both sweeps. The simulator is static client-side HTML/CSS/JavaScript with no backend or key; the notebook needs only Python. Inputs are synthetic: one payoff law read by both. Shortest run: open the notebook and choose \emph{Run all}.

\subsection*{Statement of Contribution to the UN's SDGs}
\textbf{Hugging Face:} \DemoLink\\
Through an open-access game, this project supports \textbf{SDG 4---Quality Education}~\cite{sdg4} by helping students and the public learn why honest disclosure requires cheap verification, by comparing the six rule-and-action combinations. No account or installation is needed. The resource released is the simulator; learning gains remain to be evaluated.

\subsection*{Submission milestones}
This is v1, submitted September 13, 2026, 11:00 P.M. Beijing time (UTC+8). Class review: September 14. Author response: September 16. Reviewer follow-up: September 17. Revised submission v2: September 20. Appendix~\ref{app:abstract} records each reflection and revision; row~8 is marked pending, as no peer review has been received.
